\chapter{The Present Tense}

Greek verbs carry their subject inside their ending. \PSEGreek{λύω} means \emph{I
loosen}, \PSEGreek{λύεις} \emph{you loosen}; no pronoun is needed.

\section{Present active endings}

\begin{table}[ht]
  \centering
  \begin{tabular}{lll}
    \textbf{Person} & \textbf{Singular} & \textbf{Plural} \\
    1st & \PSEGreek{λύω} & \PSEGreek{λύομεν} \\
    2nd & \PSEGreek{λύεις} & \PSEGreek{λύετε} \\
    3rd & \PSEGreek{λύει} & \PSEGreek{λύουσι} \\
  \end{tabular}
  \caption{The present active of \PSEGreek{λύω}.}
\end{table}

\section{Vocabulary dowry}

Verbs: \PSEGreek{λύω}, \emph{loosen}; \PSEGreek{γράφω}, \emph{write}; \PSEGreek{ἔχω}, \emph{have}.
Nouns: \PSEGreek{ὁ} \PSEGreek{λόγος}, \emph{word}; \PSEGreek{ὁ} \PSEGreek{ἄνθρωπος}, \emph{human being}; \PSEGreek{τὸ} \PSEGreek{βιβλίον},
\emph{book}. Ten minutes of writing these from memory will fix them better
than an hour of rereading.

\section{Reading practice}

\PSEGreek{λύομεν} \PSEGreek{τὸν} \PSEGreek{λόγον} — we loosen the word (we set the matter free).
\PSEGreek{ὁ} \PSEGreek{ἄνθρωπος} \PSEGreek{γράφει} \PSEGreek{βιβλίον} — the person writes a book.

\PSEIndex{present tense}
\PSEIndex{verbs}
